   \title{Deep Learning for Lunar Surface Feature Recognition:
          A Comparative Study of Segmentation and Detection Architectures}

   \author{G. Venturato\inst{1}
          \and
          A. Rahpeyma\inst{1}
          \and
          A.~M. Saiedi Saber\inst{1}
          \and
          G. Vagnoni\inst{1}
          }

   \institute{Physics of Data, Department of Physics and Astronomy,
              University of Padova, Via Marzolo 8, I-35131 Padova, Italy}

   \date{Received June 13, 2026; accepted}

   \abstract
   % context heading (optional)
    {Lunar reconnaissance missions generate high-resolution imagery at volumes that make manual geological mapping infeasible, driving the need for automated detection and segmentation of surface features.}
   % aims heading (mandatory)
    {This work compares four approaches to lunar geomorphological feature recognition---a classical thresholding-and-morphology baseline, semantic segmentation (SmallUNet), instance segmentation (Mask R-CNN), and single-stage object detection (YOLOv8)---on the same dataset, the same leakage-free spatial split, and the same metrics, in order to establish which formulation is preferable for which feature class and at what computational cost.}
   % methods heading (mandatory)
    {All methods are trained and evaluated on 15\,931 overlapping tiles of the Marius Hills region carrying seven semantic feature channels. A shared evaluation protocol reduces every method to per-class binary masks on a spatially disjoint validation set, scored with prevalence-robust metrics (Matthews correlation alongside F1/IoU) and per-method calibrated thresholds. Dedicated experiments quantify two methodological hazards of the dataset: the near-saturated impact-crater channel (82\% pixel prevalence) and the train--validation leakage induced by 50\%-overlapping tiles under random splitting.}
   % results heading (mandatory)
    {No method beats a trivial predict-everything baseline on the saturated crater channel (all MCC $\approx 0$ on dense tiles); on crater-sparse tiles Mask R-CNN leads (MCC 0.19) ahead of YOLOv8 (0.16). On wrinkle ridges---the class with genuine supervision---the ranking reverses: SmallUNet is best (MCC 0.45) ahead of Mask R-CNN (0.37), with YOLOv8 (0.08) and a classical Sato-filter baseline ($\approx 0$) far behind, at 1/60th of Mask R-CNN's inference cost. Identical YOLOv8 detectors retrained under a random tile split versus the spatial split quantify the protocol's impact on reported scores.}
   % conclusions heading (optional), leave it empty if necessary
     {The preferable architecture depends on the class regime: semantic segmentation wins where dense pixel supervision exists, instance detection wins for sparse countable objects, and no architecture can rescue a channel without genuine supervision. Evaluation-protocol choices (split geometry, metric, threshold) shift reported scores by more than any architectural choice studied here, and must be fixed before architectures are compared.}

    \keywords{Moon -- lunar surface features -- deep learning -- instance segmentation -- semantic segmentation -- object detection}

   \titlerunning{Deep Learning for Lunar Surface Feature Recognition}
   \authorrunning{Venturato et al.}

   \maketitle
